\documentclass[12pt]{article}
\usepackage{latexblog}

% ============================================================
% Blog Metadata
% ============================================================
\blogtitle{使用TDD开发SpringBoot应用}
\blogdate{2019-09-04}
\blogtags{开发实践, SpringBoot, 闲话编程}
\bloglang{zh}
\blogsummary{}

\begin{document}
\makeblogtitle

虽然觉得TDD没什么卵用，但实际工作中还是必须要使用TDD，这不最近就做了一个使用TDD的方式开发SpringBoot的例子。

下面阐述一下如何开发一个RestFul的GET请求，从数据库中读取数据并返回。

\section{创建Integration Test}\label{ux521bux5efaintegration-test}

第一步可以从IntegrationTest开始，即模拟真实发送一个HTTP请求，然后检验返回的Response。简单起见，第一步只校验状态值：

\begin{Shaded}
\begin{Highlighting}[]
\AttributeTok{@RunWith}\OperatorTok{(}\NormalTok{SpringRunner}\OperatorTok{.}\FunctionTok{class}\OperatorTok{)}
\AttributeTok{@SpringBootTest}\OperatorTok{(}\NormalTok{webEnvironment }\OperatorTok{=}\NormalTok{ WebEnvironment}\OperatorTok{.}\FunctionTok{RANDOM\_PORT}\OperatorTok{)}
\KeywordTok{public} \KeywordTok{class}\NormalTok{ IntegrationTest }\OperatorTok{\{}

  \AttributeTok{@Autowired}
  \KeywordTok{private}\NormalTok{ TestRestTemplate testClient}\OperatorTok{;}

  \AttributeTok{@Test}
  \KeywordTok{public} \DataTypeTok{void} \FunctionTok{should\_get\_computer\_list\_when\_call\_list\_computer\_api}\OperatorTok{()} \OperatorTok{\{}
\NormalTok{    ResponseEntity}\OperatorTok{\textless{}}\BuiltInTok{List}\OperatorTok{\textless{}}\NormalTok{ComputerDto}\OperatorTok{\textgreater{}\textgreater{}}\NormalTok{ response }\OperatorTok{=}\NormalTok{ testClient}\OperatorTok{.}\FunctionTok{exchange}\OperatorTok{(}
        \StringTok{"/computers"}\OperatorTok{,}
\NormalTok{        HttpMethod}\OperatorTok{.}\FunctionTok{GET}\OperatorTok{,}
        \KeywordTok{null}\OperatorTok{,}
        \KeywordTok{new}\NormalTok{ ParameterizedTypeReference}\OperatorTok{\textless{}}\BuiltInTok{List}\OperatorTok{\textless{}}\NormalTok{ComputerDto}\OperatorTok{\textgreater{}\textgreater{}()} \OperatorTok{\{}
        \OperatorTok{\});}
    \FunctionTok{assertEquals}\OperatorTok{(}\NormalTok{HttpStatus}\OperatorTok{.}\FunctionTok{OK}\OperatorTok{,}\NormalTok{ response}\OperatorTok{.}\FunctionTok{getStatusCode}\OperatorTok{());}
  \OperatorTok{\}}
\OperatorTok{\}}
\end{Highlighting}
\end{Shaded}

这里特意使用WebEnvironment.RANDOM\_PORT以使得Spring启动一个接近真实的Server，来测试我们的请求。当然这个测试会挂了，因为Controller都还没写呢。所以下一个先来创建Controller，但是呢，TDD通常从测试开始写起，所以来测试Controller吧。

\section{Controller Test}\label{controller-test}

测试Controller就是单元测试了，不需要测试其他的组件（比如service什么的)。

\begin{Shaded}
\begin{Highlighting}[]
\AttributeTok{@RunWith}\OperatorTok{(}\NormalTok{SpringRunner}\OperatorTok{.}\FunctionTok{class}\OperatorTok{)}
\AttributeTok{@WebMvcTest}\OperatorTok{(}\NormalTok{controllers }\OperatorTok{=}\NormalTok{ ComputerController}\OperatorTok{.}\FunctionTok{class}\OperatorTok{)}
\KeywordTok{public} \KeywordTok{class}\NormalTok{ ComputerControllerTest }\OperatorTok{\{}

  \AttributeTok{@Autowired}
  \KeywordTok{private}\NormalTok{ MockMvc mockMvc}\OperatorTok{;}

  \AttributeTok{@Test}
  \KeywordTok{public} \DataTypeTok{void} \FunctionTok{should\_get\_a\_list\_when\_get\_computers}\OperatorTok{()} \KeywordTok{throws} \BuiltInTok{Exception} \OperatorTok{\{}
\NormalTok{    mockMvc}\OperatorTok{.}\FunctionTok{perform}\OperatorTok{(}\NormalTok{MockMvcRequestBuilders}\OperatorTok{.}\FunctionTok{get}\OperatorTok{(}\StringTok{"/computers"}\OperatorTok{))}
        \OperatorTok{.}\FunctionTok{andExpect}\OperatorTok{(}\FunctionTok{status}\OperatorTok{().}\FunctionTok{isOk}\OperatorTok{());}
  \OperatorTok{\}}
\OperatorTok{\}}
\end{Highlighting}
\end{Shaded}

然后就是需要创建一个Controller，让测试可以过：

\begin{Shaded}
\begin{Highlighting}[]
\AttributeTok{@RestController}\OperatorTok{(}\StringTok{"/computers"}\OperatorTok{)}
\KeywordTok{public} \KeywordTok{class}\NormalTok{ ComputerController }\OperatorTok{\{}

  \AttributeTok{@GetMapping}
  \KeywordTok{public} \BuiltInTok{List}\OperatorTok{\textless{}}\NormalTok{ComputerDto}\OperatorTok{\textgreater{}} \FunctionTok{getComputers}\OperatorTok{()} \OperatorTok{\{}
    \ControlFlowTok{return} \KeywordTok{null}\OperatorTok{;}
  \OperatorTok{\}}
\OperatorTok{\}}
\end{Highlighting}
\end{Shaded}

到这里基本上集成测试也可以过了。所以你可以先commit一次了。然后，当然我们不能把逻辑放到Controller里面啊，我们需要一个Service来处理业务逻辑。这个Service又会从数据库中读取数据。在用到Repository之前，我们可能需要先改一下我们的controller测试，因为到目前为止并没有校验实际的字段，只是校验了返回状态码，现在可以开始校验了：

\begin{Shaded}
\begin{Highlighting}[]
  \AttributeTok{@MockBean}
  \KeywordTok{private}\NormalTok{ ComputerService computerService}\OperatorTok{;}

  \AttributeTok{@Test}
  \KeywordTok{public} \DataTypeTok{void} \FunctionTok{should\_get\_a\_list\_when\_get\_computers}\OperatorTok{()} \KeywordTok{throws} \BuiltInTok{Exception} \OperatorTok{\{}
    \FunctionTok{given}\OperatorTok{(}\NormalTok{computerService}\OperatorTok{.}\FunctionTok{getComputers}\OperatorTok{())}
        \OperatorTok{.}\FunctionTok{willReturn}\OperatorTok{(}
            \BuiltInTok{Collections}\OperatorTok{.}\FunctionTok{singletonList}\OperatorTok{(}
                \KeywordTok{new} \FunctionTok{ComputerDto}\OperatorTok{(}\DecValTok{1}\OperatorTok{,} \StringTok{"MacBook 2015"}\OperatorTok{,} \StringTok{"Haifeng Li"}\OperatorTok{,} \StringTok{"2019{-}09{-}10"}\OperatorTok{)}
            \OperatorTok{));}
\NormalTok{    mockMvc}\OperatorTok{.}\FunctionTok{perform}\OperatorTok{(}\NormalTok{MockMvcRequestBuilders}\OperatorTok{.}\FunctionTok{get}\OperatorTok{(}\StringTok{"/computers"}\OperatorTok{))}
        \OperatorTok{.}\FunctionTok{andExpect}\OperatorTok{(}\FunctionTok{status}\OperatorTok{().}\FunctionTok{isOk}\OperatorTok{())}
        \OperatorTok{.}\FunctionTok{andExpect}\OperatorTok{(}\FunctionTok{jsonPath}\OperatorTok{(}\StringTok{"$"}\OperatorTok{,} \FunctionTok{hasSize}\OperatorTok{(}\DecValTok{1}\OperatorTok{)))}
        \OperatorTok{.}\FunctionTok{andExpect}\OperatorTok{(}\FunctionTok{jsonPath}\OperatorTok{(}\StringTok{"$[0].id"}\OperatorTok{).}\FunctionTok{value}\OperatorTok{(}\DecValTok{1}\OperatorTok{))}
        \OperatorTok{.}\FunctionTok{andExpect}\OperatorTok{(}\FunctionTok{jsonPath}\OperatorTok{(}\StringTok{"$[0].type"}\OperatorTok{).}\FunctionTok{value}\OperatorTok{(}\StringTok{"MacBook 2015"}\OperatorTok{))}
        \OperatorTok{.}\FunctionTok{andExpect}\OperatorTok{(}\FunctionTok{jsonPath}\OperatorTok{(}\StringTok{"$[0].owner"}\OperatorTok{).}\FunctionTok{value}\OperatorTok{(}\StringTok{"Haifeng Li"}\OperatorTok{))}
        \OperatorTok{.}\FunctionTok{andExpect}\OperatorTok{(}\FunctionTok{jsonPath}\OperatorTok{(}\StringTok{"$[0].createTime"}\OperatorTok{).}\FunctionTok{value}\OperatorTok{(}\StringTok{"2019{-}09{-}10"}\OperatorTok{))}
        \OperatorTok{.}\FunctionTok{andDo}\OperatorTok{(}\FunctionTok{print}\OperatorTok{());}
  \OperatorTok{\}}
\end{Highlighting}
\end{Shaded}

这里我们把Service给Mock掉，因此可以控制它的行为，实际的Service就一个空函数就可以了。

\section{Service}\label{service}

这时候，可以考虑实现Service了，因为Service需要读取数据库，所以Service需要引入一个Repository来查询数据库，我们可以Mock掉Repository，来测service的逻辑：

\begin{Shaded}
\begin{Highlighting}[]
\AttributeTok{@RunWith}\OperatorTok{(}\NormalTok{MockitoJUnitRunner}\OperatorTok{.}\FunctionTok{class}\OperatorTok{)}
\KeywordTok{public} \KeywordTok{class}\NormalTok{ ComputerServiceTest }\OperatorTok{\{}

  \AttributeTok{@Mock}
  \KeywordTok{private}\NormalTok{ ComputerRepository computerRepository}\OperatorTok{;}

  \AttributeTok{@InjectMocks}
  \KeywordTok{private}\NormalTok{ ComputerService computerService}\OperatorTok{;}

  \AttributeTok{@Test}
  \KeywordTok{public} \DataTypeTok{void} \FunctionTok{should\_return\_computer\_list\_when\_get\_all\_computers}\OperatorTok{()} \KeywordTok{throws} \BuiltInTok{ParseException} \OperatorTok{\{}
\NormalTok{    ComputerEntity stored }\OperatorTok{=} \KeywordTok{new} \FunctionTok{ComputerEntity}\OperatorTok{(}\DecValTok{1}\OperatorTok{,}
        \StringTok{"MacBook 2015"}\OperatorTok{,}
        \StringTok{"Haifeng Li"}\OperatorTok{,}
        \KeywordTok{new} \BuiltInTok{SimpleDateFormat}\OperatorTok{(}\StringTok{"dd/MM/yyyy"}\OperatorTok{).}\FunctionTok{parse}\OperatorTok{(}\StringTok{"01/09/2019"}\OperatorTok{));}
    \FunctionTok{given}\OperatorTok{(}\NormalTok{computerRepository}\OperatorTok{.}\FunctionTok{findAll}\OperatorTok{()).}\FunctionTok{willReturn}\OperatorTok{(}\BuiltInTok{Collections}\OperatorTok{.}\FunctionTok{singletonList}\OperatorTok{(}\NormalTok{stored}\OperatorTok{));}

    \BuiltInTok{List}\OperatorTok{\textless{}}\NormalTok{ComputerDto}\OperatorTok{\textgreater{}}\NormalTok{ computers }\OperatorTok{=}\NormalTok{ computerService}\OperatorTok{.}\FunctionTok{getComputers}\OperatorTok{();}

    \FunctionTok{assertEquals}\OperatorTok{(}\DecValTok{1}\OperatorTok{,}\NormalTok{ computers}\OperatorTok{.}\FunctionTok{size}\OperatorTok{());}
    \FunctionTok{assertEquals}\OperatorTok{(}\DecValTok{1}\OperatorTok{,}\NormalTok{ computers}\OperatorTok{.}\FunctionTok{get}\OperatorTok{(}\DecValTok{0}\OperatorTok{).}\FunctionTok{getId}\OperatorTok{());}
    \FunctionTok{assertEquals}\OperatorTok{(}\StringTok{"MacBook 2015"}\OperatorTok{,}\NormalTok{ computers}\OperatorTok{.}\FunctionTok{get}\OperatorTok{(}\DecValTok{0}\OperatorTok{).}\FunctionTok{getType}\OperatorTok{());}
    \FunctionTok{assertEquals}\OperatorTok{(}\StringTok{"Haifeng Li"}\OperatorTok{,}\NormalTok{ computers}\OperatorTok{.}\FunctionTok{get}\OperatorTok{(}\DecValTok{0}\OperatorTok{).}\FunctionTok{getOwner}\OperatorTok{());}
    \FunctionTok{assertEquals}\OperatorTok{(}\StringTok{"2019{-}09{-}01"}\OperatorTok{,}\NormalTok{ computers}\OperatorTok{.}\FunctionTok{get}\OperatorTok{(}\DecValTok{0}\OperatorTok{).}\FunctionTok{getCreateTime}\OperatorTok{());}
  \OperatorTok{\}}
\OperatorTok{\}}
\end{Highlighting}
\end{Shaded}

同样Repository里面也就一个空函数就行了，但是这时候得把Service 的逻辑写完，让测试可以通过，这样Service的任务就完成了，其他测试也全部都可以通过。

\section{Repository}\label{repository}

最后一步就是来实现Repository了，这里需要使用DataJpaTest，用内存数据库进行测试：

\begin{Shaded}
\begin{Highlighting}[]
\AttributeTok{@RunWith}\OperatorTok{(}\NormalTok{SpringRunner}\OperatorTok{.}\FunctionTok{class}\OperatorTok{)}
\AttributeTok{@DataJpaTest}
\KeywordTok{public} \KeywordTok{class}\NormalTok{ ComputerRepositoryTest }\OperatorTok{\{}

  \AttributeTok{@Autowired}
  \KeywordTok{private}\NormalTok{ TestEntityManager entityManager}\OperatorTok{;}

  \AttributeTok{@Autowired}
  \KeywordTok{private}\NormalTok{ ComputerRepository computerRepository}\OperatorTok{;}

  \AttributeTok{@Before}
  \KeywordTok{public} \DataTypeTok{void} \FunctionTok{prepareData}\OperatorTok{()} \KeywordTok{throws} \BuiltInTok{ParseException} \OperatorTok{\{}
\NormalTok{    entityManager}\OperatorTok{.}\FunctionTok{persistAndFlush}\OperatorTok{(}\KeywordTok{new} \FunctionTok{ComputerEntity}\OperatorTok{(}\DecValTok{1}\OperatorTok{,}
        \StringTok{"MacBook 2015"}\OperatorTok{,}
        \StringTok{"Haifeng Li"}\OperatorTok{,}
        \KeywordTok{new} \BuiltInTok{SimpleDateFormat}\OperatorTok{(}\StringTok{"dd/MM/yyyy"}\OperatorTok{).}\FunctionTok{parse}\OperatorTok{(}\StringTok{"01/09/2019"}\OperatorTok{))}
    \OperatorTok{);}
\NormalTok{    entityManager}\OperatorTok{.}\FunctionTok{persistAndFlush}\OperatorTok{(}\KeywordTok{new} \FunctionTok{ComputerEntity}\OperatorTok{(}\DecValTok{2}\OperatorTok{,}
        \StringTok{"Desktop"}\OperatorTok{,}
        \KeywordTok{null}\OperatorTok{,}
        \KeywordTok{new} \BuiltInTok{SimpleDateFormat}\OperatorTok{(}\StringTok{"dd/MM/yyyy"}\OperatorTok{).}\FunctionTok{parse}\OperatorTok{(}\StringTok{"02/09/2019"}\OperatorTok{))}
    \OperatorTok{);}
  \OperatorTok{\}}

  \AttributeTok{@Test}
  \KeywordTok{public} \DataTypeTok{void} \FunctionTok{should\_return\_all\_records\_in\_db\_when\_find\_all}\OperatorTok{()} \OperatorTok{\{}
    \BuiltInTok{List}\OperatorTok{\textless{}}\NormalTok{ComputerEntity}\OperatorTok{\textgreater{}}\NormalTok{ entities }\OperatorTok{=}\NormalTok{ computerRepository}\OperatorTok{.}\FunctionTok{findAll}\OperatorTok{();}
    \FunctionTok{assertEquals}\OperatorTok{(}\DecValTok{2}\OperatorTok{,}\NormalTok{ entities}\OperatorTok{.}\FunctionTok{size}\OperatorTok{());}
  \OperatorTok{\}}
\OperatorTok{\}}
\end{Highlighting}
\end{Shaded}

因为Spring JPA你只需要写一堆interface，测试这里的逻辑还是十分有必要的。所以到这一步为止，基本上程序的功能就已经实现了，唯一需要做的就是改一下配置来连接到真是的数据库。整个代码可以在我的\href{https://github.com/soleverlee/computer-inventory.git}{Github}上面找到。


\end{document}
